% ============================================================================
%  The Self-Consistency Ceiling Is Not a Model Constant:
%  Modal Answers Are a Model x Elicitation Property
%
%  STYLE NOTE / HOW TO SWAP IN A CONFERENCE STYLE
%  ----------------------------------------------
%  This file compiles standalone with the stock `article` class so that it
%  builds anywhere (`pdflatex main && bibtex main && pdflatex main x2`) with no
%  downloaded style files.
%
%  To switch to a conference format:
%    * NeurIPS : download neurips_2026.sty, then replace the
%                \documentclass line with \documentclass{article} and add
%                \usepackage[preprint]{neurips_2026}
%    * ICLR    : download iclr2027_conference.sty + .bst, replace with
%                \usepackage{iclr2027_conference,times}
%    * ICML    : download icml2026.sty, use \usepackage[accepted]{icml2026}
%    * ACL/EMNLP : download acl.sty + acl_natbib.bst, use
%                \documentclass[11pt]{article} \usepackage[]{acl}
%  In every case, delete the \geometry{...} line below (the conference style
%  sets the page geometry itself) and keep everything from \begin{document} on.
% ============================================================================
\documentclass[11pt]{article}

\usepackage[margin=1in]{geometry}
\usepackage[utf8]{inputenc}
\usepackage[T1]{fontenc}
\usepackage{times}
\usepackage{amsmath,amssymb,amsthm}
\usepackage{graphicx}
\usepackage{booktabs}
\usepackage{multirow}
\usepackage{xcolor}
\usepackage{url}
\usepackage[numbers,sort&compress]{natbib}
\usepackage[colorlinks=true,linkcolor=blue,citecolor=blue,urlcolor=blue]{hyperref}
\usepackage{algorithm}
\usepackage{algpseudocode}

% --- \todo marker (self-contained; no todonotes dependency) -----------------
\newcommand{\todo}[1]{\textcolor{red}{\textbf{[TODO: #1]}}}
\newcommand{\placeholder}[1]{\textcolor{gray}{\textit{#1}}}

% --- notation ---------------------------------------------------------------
\newcommand{\pmode}{\pi_{\mathrm{mode}}}
\newcommand{\phat}{\hat{p}}
\newcommand{\rdis}{\rho_{\mathrm{disatt}}}

\title{\bfseries The Self-Consistency Ceiling Is Not a Model Constant:\\
Modal Answers Are a Model\,$\times$\,Elicitation Property}

\author{
  Anonymous Author(s) \\
  Affiliation \\
  \texttt{email@domain}
}

\date{}

\begin{document}
\maketitle

% ============================================================================
\begin{abstract}
% ============================================================================
Recent work bounds self-consistency by a \emph{modal ceiling} $\pmode$: the
fraction of items whose plurality answer is correct. Coverage keeps climbing;
selection plateaus. That wedge is the \emph{identifiability gap}. The
literature treats $\pmode$ as a property of a model on a benchmark.
Contemporaneous analyses show per-item difficulty is already substantially
prompt-attributable on log-prob-scored benchmarks
\citep{brittlebench2026,tee2026} and that unstable difficulty labels corrupt
downstream filters \citep{hardunreached2026}. We ask whether prompt sensitivity
reaches the \textbf{mode} in \textbf{sampled chain-of-thought} on
free-generation math---where inference variance does not vanish---and whether
spreading a fixed token budget across configurations can \emph{raise} the
plateau that i.i.d.\ resampling structurally cannot move.

On a fully crossed grid over items, three Qwen sizes, a Llama cross-family
control, and three core semantics-preserving configurations (c0--c2;
independent-seed replicates), we measure $\pmode$ per configuration, mode
reordering on small-margin items, and configuration-diversified voting (CDV),
which repartitions an existing corpus at zero extra GPU cost. A crossed
binomial GLMM provides supporting variance decomposition
\citep{tee2026,prompteval2024}. Primary predictions: $\pmode$ varies by ${>}3$
points across core configurations (P4); plug-in $\pmode$(CDV) exceeds
$\pmode$(c0) at matched token budget while temperature-diversified SC (B4)
does not (P5/F3). Mode transitions split into corrective vs.\ destructive
components (Figure~F6).
Our results imply that reported reasoning boundaries are properties of a
model--prompt pair, not of a model alone, and that configuration diversity is
an orthogonal axis to resampling depth for moving the modal ceiling.
\end{abstract}

% ============================================================================
\section{Introduction}
% ============================================================================

Progress in language-model reasoning is increasingly mediated by how a model
\emph{selects} an answer under sampling. Compute-optimal test-time scaling
bins items by estimated difficulty~\citep{snell2024scaling}; adaptive
self-consistency stops when a vote looks decided~\citep{aggarwal2023adaptive,li2024esc};
and the analyses that explain why extra samples stop helping bound selection by
the \emph{modal-hit rate} $\pmode$---the fraction of items whose plurality
answer is correct~\citep{brown2024monkeys,blendasc2025,certifiedsc2025}.
More i.i.d.\ samples estimate the mode more precisely; they cannot change which
class is modal. If elicitation configuration changes the mode, then $\pmode$ is
not a model constant and there is headroom that resampling structurally cannot
access.

That ``elicitation matters'' is no longer an open question at the item level.
\citet{brittlebench2026} decompose benchmark variance into task difficulty and
semantics-preserving prompt sensitivity and find the latter can account for
roughly half of total variance on multiple-choice suites scored deterministically.
\citet{tee2026} fit crossed random effects with an explicit item\,$\times$\,prompt
interaction by REML. \citet{hardunreached2026} show sampling-derived hardness
labels are unstable and corrupt RL filters and curricula. What these works do
\textbf{not} examine is whether sensitivity reaches the \textbf{mode} on
\textbf{sampled} chain-of-thought for free-generation math, where binomial noise
is load-bearing, or whether spreading budget across configurations can raise
$\pmode$ rather than merely approach it faster.

We therefore ask three questions. \textbf{(Q1)} Does $\pmode$ move across
semantics-preserving elicitation configurations for the same item and model?
\textbf{(Q2)} If so, is the movement explained by \emph{reordering} of top
answer classes on small-margin items? \textbf{(Q3)} Can
\emph{configuration-diversified voting} (CDV)---pooling a fixed token budget
across configurations---raise the plateau above single-configuration
self-consistency, while temperature-diversified sampling at matched budget does
not? Supporting measurement (Q4): how much item\,$\times$\,configuration
structure exists in a crossed binomial GLMM, and do difficulty labels transfer?

Answering these requires a design that most prior work does not have. A fully
crossed item\,$\times$\,model\,$\times$\,configuration grid is necessary to
attribute variance. Independent-seed replicates within a single configuration
are necessary because $\phat = k/N$ is a noisy estimate and any two estimates
will disagree somewhat for purely binomial reasons; without a measured noise
floor, an ``interaction'' is unfalsifiable. And a budget-matched comparison in
\emph{output tokens} rather than sample count is necessary because
configurations induce chains of different lengths. We build all three in.

\paragraph{Contributions.}
\begin{enumerate}\itemsep2pt
  \item \textbf{$\pmode$ is configuration-relative.} We measure the modal-hit
    rate per elicitation configuration on MATH-500 and GSM-Symbolic with sampled
    CoT, across a size ladder and two families (Figure~F4).
  \item \textbf{A mechanism: mode reordering.} Configuration changes reorder top
    answer classes on small-margin items; we separate corrective from
    destructive transitions (Figure~F6).
  \item \textbf{CDV raises the plateau.} Configuration-diversified voting spreads
    a fixed token budget across configurations and is pre-registered to raise
    plug-in $\pmode$ above single-configuration self-consistency, unlike
    temperature-diversified sampling (Figure~F5; primary test P5).
  \item \textbf{Supporting measurement.} Crossed binomial GLMM variance
    decomposition \citep{tee2026}, disattenuated transfer, and hard-subset
    Jaccard against a seed null (Figures~F1--F3), calibrated by a
    parametric-bootstrap null on synthetic data.
  \item \textbf{Applied consequence.} Transferred-difficulty allocation degrades
    when difficulty is estimated in one configuration and spent in another; curriculum
    RL discussed as implication only \citep{hardunreached2026}.
  \item \textbf{Pre-registered predictions} with individually decisive
    falsification conditions (F3 primary); stated negative-result paths.
\end{enumerate}

\paragraph{Scope and honesty about what is new.} We do not claim to be the first
to show that prompts affect per-item outcomes: BrittleBench, TEE, PromptEval,
and FormatSpread establish prompt sensitivity and variance decomposition in
adjacent settings \citep{brittlebench2026,tee2026,prompteval2024,sclar2024formatspread,mizrahi2024multiprompt}.
Our claim is that sensitivity reaches the \emph{mode}---the object that caps
self-consistency---on \emph{sampled} reasoning chains, and that this is
\emph{exploitable} without additional generation.

% ============================================================================
\section{Related Work}
\label{sec:related}
% ============================================================================

\subsection{Chain-of-thought and its structural variants}

Chain-of-thought prompting~\citep{wei2022cot} and its zero-shot
form~\citep{kojima2022zeroshot} established that eliciting intermediate steps
improves reasoning accuracy, and self-consistency~\citep{wang2023selfconsistency}
established that sampling several chains and taking a plurality vote improves
it further. A family of structural elaborations followed:
least-to-most decomposition~\citep{zhou2023leasttomost}, tree-~\citep{yao2023tot}
and graph-structured~\citep{besta2024got} search over partial thoughts, and
program-mediated reasoning that offloads arithmetic to an
interpreter~\citep{chen2023pot,gao2023pal}. Later work tempered the picture:
\citet{sprague2025tocot} find that CoT's benefit is concentrated on
mathematical and symbolic tasks and is small or negative elsewhere, and much of
the apparent advantage of elaborate search structures narrows once sampling
budget is matched. \textbf{What remains unresolved:} the community still lacks
a principled account of \emph{which} items benefit from which structure, and
almost all such accounts are stated in terms of a per-item difficulty that, we
argue here, is itself unstable.

\subsection{Test-time compute scaling and its ceilings}

Our benchmarks are \textbf{MATH-500}~\citep{math500} (headline) and
GSM-Symbolic~\citep{mirzadeh2025gsmsymbolic}, with GSM8K~\citep{cobbe2021gsm8k}
as a secondary replicate. Verifier-reranked sampling~\citep{cobbe2021gsm8k}, process
supervision~\citep{lightman2024verify,uesato2022process,wang2024mathshepherd},
and compute-optimal allocation between sequential revision and parallel
sampling~\citep{snell2024scaling,wu2024inference} form the backbone of
test-time scaling; \citet{brown2024monkeys} showed that coverage---the
probability that at least one of $n$ samples is correct---keeps rising over
orders of magnitude of $n$, while plurality voting and learned reward models
plateau. \citet{muennighoff2025s1} showed that a very simple budget-forcing
scheme captures much of the benefit. Surveys now organise the
area along what/how/where/how-well axes~\citep{zhang2025ttssurvey} and by
subproblem topology~\citep{subproblemsurvey2025}, and a large matched-condition
study finds that no single strategy dominates across models and difficulty
regimes~\citep{artofscaling2025}.

The most directly relevant recent development is the explanation of the
plateau. \citet{brown2024monkeys} showed coverage keeps rising while plurality
voting plateaus; subsequent work formalises a \emph{modal ceiling}:
plurality voting converges to the mode, so selection accuracy saturates at
$\pmode$; the gap to rising coverage is the \emph{identifiability gap}
\citep{blendasc2025,certifiedsc2025,samplecomplexity2025}.
\textbf{What remains unresolved:}
all of this conditions on a per-item answer distribution under a \emph{fixed}
elicitation setting. \citet{bay2026modal} estimate run-to-run variance within a
configuration ($\rho_w \approx 0$) but never vary the prompt; whether $\pmode$
is configuration-relative is the step we test.

\subsection{Prompt sensitivity and variance decomposition}
\label{sec:related-prompt}

Three 2026 lines establish that per-item outcomes are prompt-attributable and
that evaluation pipelines must model crossed facets. \textbf{BrittleBench}
\citep{brittlebench2026} decomposes binary correctness variance into task
difficulty and semantics-preserving prompt sensitivity on an item\,$\times$\,perturbation
matrix; inference is deterministic (log-prob scoring), so the inference-variance
term vanishes. \textbf{TEE} \citep{tee2026} fits crossed random effects---including
item\,$\times$\,prompt---by REML and uses decision studies on aggregate benchmark
scores. \textbf{PromptEval} \citep{prompteval2024} estimates performance
distributions across prompt templates with an IRT-style model on the joint
(item, template) matrix. Earlier aggregate work
\citep{sclar2024formatspread,mizrahi2024multiprompt,alzahrani2024leaderboard}
showed format and template choice move rankings by large margins.

\textbf{What remains unresolved} is the intersection with \emph{sampled}
chain-of-thought on free-generation math and with the \emph{mode}. BrittleBench
does not report $\pmode$, reorder rates, or CDV; TEE's estimand is
$\mathrm{Var}(\hat\theta)$ at benchmark scale---item\,$\times$\,prompt is a
nuisance that averages away in aggregate CIs, whereas adaptive test-time compute,
hard-subset curation, and IRT consumers use \emph{per-item} parameters one at a
time. \citet{facetnoise2026,cyclicjudge2026} show that under G-theory facet decompositions,
repeats are the cheap facet and paraphrases/models buy more reliability per
query---the allocation-theoretic sibling of ``i.i.d.\ resampling cannot move the
plateau,'' but on brand sentiment without modes or reasoning. \citet{hardunreached2026}
establish unstable hardness labels along the stochastic-vs-deterministic axis;
we study the \emph{configuration} axis and whether it reaches the selection
bottleneck. Adjacent aggregation devices---paraphrase
voting~\citep{paraconsistency2024,zhou2024paraphrase}---are baselines, not the
measurement target.

\subsection{Long chain-of-thought, RL for reasoning, and the elicitation debate}

Reasoning-specialised models trained with reinforcement learning on verifiable
rewards~\citep{deepseekr1_2025,shao2024deepseekmath} produce long
chains-of-thought and large gains, and distillation transfers much of the
behaviour into small models with remarkably little
data~\citep{muennighoff2025s1,ye2025limo}. A pointed debate followed about what
RL actually contributes. \citet{yue2025rlvr} argue via pass@$k$ analysis that
RLVR largely sharpens the base model's existing distribution rather than
extending its reach, and \citet{diversitycollapse2025} formalise a related
bias--variance decomposition of pass@$k$, showing that supervised finetuning
improves pass@1 while collapsing generation diversity and degrading pass@$k$,
and that weight interpolation partially recovers both. Theoretically,
self-improvement has been framed as \emph{sharpening} a fixed base
policy~\citep{huang2025sharpening}, and inference-time methods have been
recast as sharpening the answer marginal
directly~\citep{marginalsharpening2026}. \textbf{What remains unresolved:} the
sharpening framing implies that the ordering of candidate answers is doing all
the work, but the stability of that ordering---under scale, under
post-training, and under elicitation---has not been measured. Our
mode-reordering analysis (\S\ref{sec:method-reorder}) is the missing
measurement for the elicitation axis.

\subsection{Faithfulness, premature commitment, and what the chain is for}

A parallel literature asks whether the chain reflects the computation.
\citet{turpin2023unfaithful} show that models silently follow biasing features
in the context while producing unrelated rationales;
\citet{lanham2023measuring} introduce the truncation and corruption probes that
have become standard; and \citet{lyu2023faithful} propose constrained
formulations that are faithful by construction. Latent-space variants ask
whether the tokens are needed at
all~\citep{hao2024coconut,pfau2024filler,goyal2024pause}. Recent work sharpens
the commitment question: \citet{prematureconfidence2026} probe at evenly spaced
truncation points, define a premature-confidence score from the resulting
confidence trajectory, show it predicts logical flaws, and mitigate it with a
shaped RL objective; \citet{precotprobe2026} show that a linear probe on
pre-CoT activations predicts the final answer with high AUC and that the
direction is causal under steering. \citet{matcha2025} show the converse
decoupling---answers that stay correct while the rationale collapses under
imperceptible input perturbation. \textbf{What remains unresolved:} these
results establish that a model often has an answer in mind early, but they do
not connect that to the \emph{selection} bottleneck at the heart of test-time
scaling. If the answer is fixed early, then what a configuration change
perturbs is that early commitment, which is precisely why it can move a mode
that resampling cannot.

\subsection{Efficiency, overthinking, and adaptive compute}

Reasoning models overspend on easy items~\citep{overthinking2024} and
under-explore on hard ones~\citep{underthinking2025}; surveys catalogue the
resulting efficiency
literature~\citep{stopoverthinking2025}. Length-controlled policy
optimisation~\citep{l1_2025} and budget-forcing~\citep{muennighoff2025s1} give
explicit control over chain length. On the parallel axis,
adaptive-consistency~\citep{aggarwal2023adaptive} and early-stopping
self-consistency~\citep{li2024esc} halt when the vote is decided;
Certaindex/Dynasor~\citep{fu2024certaindex} generalise this to an
algorithm-agnostic ``answer stabilisation'' signal used for scheduling;
Blend-ASC~\citep{blendasc2025} allocates samples across questions to optimise
sample efficiency; \citet{certifiedsc2025} provide sequential certificates; and
\citet{adaptivettc2026} study difficulty-conditioned budget allocation directly.
\textbf{What remains unresolved:} every one of these estimates a per-item
quantity in one setting and consumes it in another, and none validates that the
estimate transfers. This is the applied consequence we measure.

\subsection{Robustness, perturbation, and the reasoning-versus-memorisation debate}

\citet{mirzadeh2025gsmsymbolic} introduced templated variants of GSM8K and
reported both high variance across instantiations of the same template and a
large drop when an irrelevant clause is added (GSM-NoOp), arguing for
pattern-matching over formal reasoning. Contamination
audits~\citep{zhang2024gsm1k}, counterfactual
evaluation~\citep{wu2024reciting}, premise-order
effects~\citep{chen2024premise} and the reversal
curse~\citep{berglund2024reversal} point the same direction, while
compositionality analyses~\citep{dziri2023faithfate} give a mechanistic account
of multi-step failure. The debate is live: puzzle-complexity critiques and
their rebuttals continue, and 2026 replications of GSM-NoOp on frontier models
suggest that a substantial part of the original effect reflects models
reasonably treating ambiguous clauses as informative rather than failing to
reason \todo{cite the 2026 audited replication once a peer-reviewed version
exists; currently only a non-archival write-up, marked unverified}. Follow-up
work finds that program-mediated reasoning does not repair GSM-Symbolic
robustness~\citep{reasoningcodeboth2026}. \textbf{What remains unresolved:}
these studies compare accuracy across perturbations of the \emph{problem}. Ours
holds the problem fixed and perturbs only the \emph{elicitation}, which
isolates a different and, we argue, more troubling source of instability,
because it cannot be blamed on the item.

\subsection{Self-correction, verification, and aggregation beyond frequency}

Iterative refinement with self-generated feedback~\citep{madaan2023selfrefine}
and multi-agent debate~\citep{du2024debate} were early attempts to convert extra
inference into extra accuracy without a verifier, but intrinsic self-correction
without external feedback does not reliably improve
reasoning~\citep{huang2024cannotselfcorrect,kamoi2024whencan}, which is part of
why the field turned to verifiers and to better aggregation. Beyond plain
plurality voting there is now a dense cluster: universal
self-consistency~\citep{chen2023universal}, program-verified
voting~\citep{prove2024}, semantic weighting of reasoning
paths~\citep{weightedreasoning2024}, confidence-informed voting and its
efficient variant~\citep{veccisc2026}, learned rerankers over answer-level
features~\citep{risc2026}, radial consensus scoring for
best-of-$N$~\citep{rcs2026}, and best-of-majority
hybrids~\citep{bestofmajority2025}.
Diversity-side interventions are covered in
\S\ref{sec:related-prompt}. \textbf{What remains unresolved:}
this literature searches for a better \emph{selector} given a fixed sample
pool. We ask a prior question---what determines the pool's mode in the first
place---and show that the answer includes a factor the selector literature
treats as fixed. The strongest of these methods are our baselines.

\subsection{Uncertainty, calibration, and evaluation methodology}

Confidence signals for reasoning include self-evaluated
correctness~\citep{kadavath2022know}, semantic
entropy~\citep{kuhn2023semantic,farquhar2024semantic}, verbalised
confidence~\citep{lin2022teaching,tian2023just}, and vote share itself. On the
evaluation side, the field has become considerably more careful:
\citet{miller2024errorbars} argues for routine interval reporting,
\citet{hochlehnert2025sober} document how much of recent reported progress on
small reasoning benchmarks is seed variance, \citet{dontpassk2025} propose
Bayesian posterior reporting in place of pass@$k$,
\citet{coverattau2025} show that pass@$k$ is biased toward low-reliability
``lucky hits'' and propose a reliability-controlled alternative, and
\citet{falsepositives2025} show that final-answer matching admits false
positives at a rate that materially lowers the true inference-scaling ceiling.
Psychometric approaches fit item difficulty explicitly and use it for adaptive
testing and routing~\citep{atlas2025,irtnet2025,irtrouter2025,psnirt2026}, and
\citet{kim2025correlated} show that errors are strongly correlated across
models, capping ensembles---an effect formalised as a co-failure
ceiling~\citep{cofailure2026}. \textbf{What remains unresolved:} the
psychometric line assumes a unidimensional item difficulty and validates it
largely on non-reasoning benchmarks. We test that assumption where it is most
consequential, and we adopt the evaluation hygiene this literature demands
throughout.

% ============================================================================
\section{Method}
\label{sec:method}
% ============================================================================

Our design is deliberately austere: a single sampling corpus, from which every
quantity in the paper is computed offline. This is what makes the study
reproducible on a free-tier GPU allocation, and it means the intervention we
propose is evaluated without generating a single additional token.

\subsection{Setup and notation}

Let $\mathcal{I}$ be a set of items, $\mathcal{M}$ a set of models,
$\mathcal{C}$ a set of elicitation configurations, and $\mathcal{S}$ a set of
sampling seeds. For each cell $(i,m,c,s)$ we draw $N$ chains-of-thought,
extract a final answer from each, canonicalise it into an equivalence class,
and record whether it is correct. Write $k_{imcs}$ for the number of correct
samples and
\begin{equation}
\phat_{imcs} \;=\; \frac{k_{imcs}}{N},
\qquad
z_{imcs} \;=\; \log\frac{k_{imcs}+\tfrac12}{\,N-k_{imcs}+\tfrac12\,},
\end{equation}
where $z$ is the Haldane--Anscombe-corrected empirical logit, whose sampling
variance is approximately $v_{imcs} = (k+\tfrac12)^{-1} + (N-k+\tfrac12)^{-1}$.
We work on the logit scale because $\phat$ is a bounded proportion whose
variance depends on its mean, which would otherwise contaminate a variance
decomposition.

An \emph{elicitation configuration} perturbs everything except mathematical
content: system prompt, few-shot order (at fixed shot count), and related axes.
Three \textbf{core} configurations $\{c_0,c_1,c_2\}$ enter the crossed
grid and CDV; parse-variant ($c_3$), few-shot ($c_4$/$c_{4a}$/$c_{4b}$), and
temperature ($c_5$) are reported on separate axes. Paraphrase ($c_6$) supports baseline
B5 only.

\paragraph{Modal ceiling (self-contained).} For item $i$, model $m$,
configuration $c$, let $a^\star$ be the plurality answer class over $N$
samples. The \emph{modal-hit rate} is
$\pmode(m,c)=\frac{1}{|\mathcal I|}\sum_i \mathbf 1[a^\star\text{ correct}]$.
The \emph{identifiability gap} is $\mathrm{pass@}N-\pmode$. We cite
\citet{brown2024monkeys} for the coverage--selection distinction and
\citet{bay2026modal} for the observation that within-configuration run-to-run
variance is negligible; we measure $\pmode$ ourselves.

\subsection{Supporting measurement: crossed binomial GLMM}
\label{sec:method-variance}

Following \citet{tee2026,prompteval2024}, we fit
\begin{equation}
k_{imc} \sim \mathrm{Binomial}(N, p_{imc}), \quad
\mathrm{logit}(p_{imc}) = \mu + a_i + b_m + g_c
  + (ab)_{im} + (ag)_{ic} + (bg)_{mc},
\end{equation}
with crossed random effects by REML. The item\,$\times$\,configuration
component $(ag)_{ic}$ is supporting evidence, calibrated against a
parametric-bootstrap null (no interaction, marginal $p_i$ only). A legacy
method-of-moments decomposition on Haldane logits is appendix robustness only.

\subsection{Disattenuated difficulty transfer}
\label{sec:method-transfer}

A raw correlation between $\phat$ under two configurations understates
agreement, because both estimates are noisy. The seed replicates let us remove
that attenuation honestly. We first compute the reliability $r_{mm}$ as the
mean pairwise rank correlation of $\phat$ across independent seeds within the
reference configuration; this is the correlation that measurement noise alone
permits. We then report, for each configuration pair, the raw rank correlation
and the disattenuated value $\rdis = \rho_{\mathrm{raw}} / r_{mm}$. The
comparison of interest is $\rdis$ against $1$: a value near one means
difficulty transfers and our claim is refuted; a value well below one means the
disagreement is structural rather than statistical. The same construction gives
transfer across model sizes and across model families, which calibrates the
configuration effect against effects the field already accepts as real.

We complement this with a set-level statistic that speaks directly to practice.
Taking the bottom quartile of $\phat$ as the ``hard subset'', we compute its
Jaccard overlap across configuration pairs and, as the null, across seed pairs.
The difference is the excess instability: the fraction of a curated hard subset
that is an artefact of the configuration in which it was curated.

\subsection{Mechanism: mode reordering}
\label{sec:method-reorder}

Self-consistency returns the modal answer class; its ceiling is $\pmode$
(\S\ref{sec:method}). We classify each item's mode transition across
configuration pairs as benign, corrective, or destructive, and predict
concentration on small top-two margins.

\subsection{Intervention: configuration-diversified voting}
\label{sec:method-cdv}

Configuration-diversified voting (CDV) spends a fixed budget across $C$
configurations instead of concentrating it in one, and takes the plurality over
the pooled samples. An adaptive variant first spends a small probe budget per
configuration and commits immediately when all configurations agree, spending
the remainder only on disagreement. Both are computed by re-partitioning the
existing corpus, so they cost nothing beyond the sampling already performed.

\begin{algorithm}[t]
\caption{Configuration-Diversified Voting (CDV)}
\begin{algorithmic}[1]
\Require item $i$, model $m$, token budget $B$, configuration set $\mathcal{C}_{\mathrm{use}}$
\State $b \gets B / |\mathcal{C}_{\mathrm{use}}|$ \Comment{split by \emph{tokens}, not samples}
\State $P \gets \emptyset$
\For{$c \in \mathcal{C}_{\mathrm{use}}$}
  \State draw chains from cell $(i,m,c)$ without replacement until their cumulative output tokens reach $b$; add their answer classes to $P$
\EndFor
\State \Return $\arg\max_{a} \#\{a \in P\}$
\end{algorithmic}
\end{algorithm}

All comparisons are matched on \emph{total generated output tokens}, not on
sample count, because configurations induce chains of different lengths;
reporting accuracy against $N$ would silently advantage whichever configuration
is most verbose.

The comparison that carries the scientific weight is CDV against
single-configuration self-consistency and against temperature-diversified SC
(B4) at matched tokens. \textbf{Primary test (P5/F3):} plug-in $\pmode$ from
$N$ samples at $c_0$ vs.\ $\pmode$ from pooled $C_{\mathrm{use}}\times N$
samples at CDV ($C_{\mathrm{use}}=3$ by default). Maj@$n$ curves are secondary; the small-budget overlap
($\leq 24$ tokens) is labelled underpowered for plateau claims. We also report
oracle best single configuration (B12) and decompose CDV gains into new modal
mass vs.\ configuration selection.

% ============================================================================
\section{Experimental Setup}
\label{sec:setup}
% ============================================================================

Our size ladder is the Qwen2.5 instruct family at 1.5B, 3B and
7B~\citep{qwen25,qwen2}, with a Llama-3 cross-family control, evaluated on
MATH-500~\citep{math500} (headline), GSM-Symbolic~\citep{mirzadeh2025gsmsymbolic},
and GSM8K~\citep{cobbe2021gsm8k} (secondary replicate).

\placeholder{The remainder of this section is finalised once the harness
reports T1. It must state: exact HuggingFace model IDs and revisions; exact
dataset IDs, splits and the persisted item-index file; the three core configuration
definitions verbatim; sampling hyperparameters; the answer-canonicalisation
procedure; the seed protocol; and measured wall-clock and token counts. All
values are specified normatively in \texttt{docs/METHOD\_SPEC.md}; this section
reports what was actually run, including any contingency that fired.}

\todo{Insert Table T1 (\texttt{tab1\_setup.csv}) once generated.}

\paragraph{Pre-registered predictions.}
Primary cell: Qwen2.5-3B $\times$ MATH-500 unless noted.
\begin{center}
\small
\begin{tabular}{@{}llp{7.5cm}@{}}
\toprule
ID & Tier & Prediction \\
\midrule
P4 & Primary & $\pmode$ varies by ${>}3$ accuracy points across core configurations. \\
P5 & Primary & Plug-in $\pmode$(CDV) $>$ $\pmode$($c_0$) at matched tokens; B4 does not exceed $\pmode$($c_0$). \\
P1 & Supporting & GLMM item$\times$config share exceeds parametric-bootstrap null 95th percentile. \\
P2 & Supporting & $\rdis$ (or GLMM latent correlation) for config-pairs ${<}0.85$; $r_{mm}$ by difficulty quintile. \\
P3 & Supporting & $J_{\mathrm{config}} \leq J_{\mathrm{seed}} - 0.10$ (interaction permutation $p{<}0.05$). \\
P6 & Applied & Transferred-difficulty allocation (B9) loses to configuration-aware allocation. \\
P7 & Gate & Item$\times$config share differs ${<}3$ pp fp16 vs.\ 4-bit; $\rdis$ differs ${<}0.05$ on 1.5B/3B. \\
\bottomrule
\end{tabular}
\end{center}
Falsification (individually decisive): \textbf{F3} plug-in $\pmode$(CDV)$\leq\pmode$($c_0$);
\textbf{F4} $\pmode$ range $\leq 3$ pp; \textbf{F1} GLMM share $\leq 5\%$ and below null;
\textbf{F2} $J_{\mathrm{config}} \geq J_{\mathrm{seed}} - 0.02$ (relative, not $J\geq 0.90$).
Mark each supported/refuted after analysis.

% ============================================================================
\section{Results}
\label{sec:results}
% ============================================================================

\placeholder{All content below is a placeholder tied to the artefact names
defined in \texttt{docs/METHOD\_SPEC.md} \S9. Do not write prose here until the
corresponding artefact exists.}

\subsection{Modal ceilings across configurations (Q1)}
\todo{Figure F4 (\texttt{fig4\_ceiling\_vs\_coverage.pdf}), Table T2. Lead Results. State P4.}

\subsection{Configuration-diversified voting (Q3)}
\todo{Figure F5, Table T5. Primary P5/F3: plug-in $\pmode$ comparison. B12 and gain decomposition.}

\subsection{Mode reordering mechanism (Q2)}
\todo{Figure F6. Corrective vs.\ destructive rates.}

\subsection{Supporting measurement: variance and transfer (Q4)}
\todo{Figures F1--F3, Tables T3--T4. GLMM primary; moment decomposition in appendix.}

\subsection{Downstream cost of transferred difficulty}
\todo{Table T6. B9 vs.\ configuration-aware allocation. P6. Curriculum RL as Discussion implication only.}

% ============================================================================
\section{Analysis}
\label{sec:analysis}
% ============================================================================

\placeholder{Placeholders tied to named artefacts; write only against real
numbers.}

\paragraph{Which items are susceptible?}
\todo{Relate reordering to the top-two answer margin. Show the predicted
concentration of reorderings on small-margin items, or report its absence.}

\paragraph{Mechanism: corrective versus destructive reordering.}
\todo{Figure F6 (\texttt{fig6\_reordering.pdf}). Report the corrective and
destructive rates by model size and discuss whether the net is positive and why.}

\paragraph{Does scale reduce the interaction?}
\todo{Compare the item$\times$configuration share across 1.5B / 3B / 7B and
across families. This is the reviewer's ``does it scale?'' question; answer it
with the trend, not with a single large model.}

\paragraph{Distribution shift.}
\todo{Figure F7 (\texttt{fig7\_shift.pdf}). Does symbolic perturbation amplify
the item$\times$configuration interaction relative to the original benchmark?}

\paragraph{Relationship to the modal-ceiling account.}
\todo{Quantify how much of the identifiability gap of \citet{bay2026modal} is
closed by configuration diversification, and how much remains for a verifier.}

\paragraph{Reconciling with unidimensional IRT.}
\todo{Report the degradation in unidimensional IRT fit when configuration is
added as a crossed factor, and the variance share such models absorb as noise.}

% ============================================================================
\section{Limitations}
% ============================================================================

\begin{itemize}\itemsep2pt
  \item \textbf{Model scale.} Every model we study runs on one or two 16\,GB
    accelerators. We address extrapolation with a size ladder and two families
    rather than a single large model, but we cannot rule out that the
    item\,$\times$\,configuration interaction shrinks at frontier scale. The
    trend across our ladder is the evidence we can offer, and we report it
    whichever way it points.
  \item \textbf{Configuration space is a sample, not a census.} Three core
    configurations along semantics-preserving axes do not span elicitation settings. Estimates
    are \textbf{conditional on our chosen configurations}, not lower bounds on a
    broader population.
  \item \textbf{Semantics preservation is argued, not proven.} We restrict the
    core configurations to axes that do not touch the problem statement, and we
    report the paraphrase axis separately for exactly this reason. A residual
    worry remains that, for instance, a persona system prompt shifts the task
    in some subtle way.
  \item \textbf{Finite $N$.} With $N{=}24$ per cell, $\phat$ has non-trivial
    binomial error. The seed-replicate arm measures and corrects for this, but
    the correction is only as good as the replicate estimate, and components
    near zero are subject to clamping bias.
  \item \textbf{Benchmark breadth.} Headline results are on MATH-500 and
    GSM-Symbolic; GSM8K is a secondary replicate (saturation and contamination).
  \item \textbf{Contamination.} GSM8K contamination is item-specific; we lead
    with MATH-500 and GSM-Symbolic and state the memorisation\,$\times$\,configuration
    risk plainly.
  \item \textbf{Answer canonicalisation.} Equivalence classes are defined by a
    symbolic checker. Its errors propagate into the modal-hit rate. We report
    extraction-failure rates and flag any cell above five percent.
  \item \textbf{Numerical precision.} Weight quantisation is itself a
    perturbation of the answer distribution, and this paper is about how
    answer distributions respond to perturbation, so the two cannot be
    casually separated. Specifically, quantisation raises perplexity, which
    flattens next-token distributions, which shrinks the top-two answer-class
    margin --- the very quantity our mechanism analysis turns on --- and so
    would be expected to \emph{inflate} the interaction we report. We
    therefore run the ladder at \texttt{fp16} and treat reduced precision as
    a measured quantity rather than an assumption: the precision control
    (Figure F8, \texttt{fig8\_precision\_control.pdf}; Table T7,
    \texttt{tab7\_precision\_control.csv}) repeats the full decomposition at four-bit on the
    two smaller rungs, where quantisation damage is relatively largest and
    the extrapolation to the larger rung therefore runs in the favourable
    direction. \todo{After the run: state whether the ladder was executed at
    \texttt{fp16} or, if the backend fallback fired, uniformly at four bits;
    in the latter case cite F8 as the licence and repeat the margin-invariance
    result here.}
\end{itemize}

% ============================================================================
\section{Conclusion}
% ============================================================================

\placeholder{Write after Results. The conclusion should state, in one
paragraph, whether $\pmode$ was configuration-relative; whether CDV raised the
modal plateau above single-configuration self-consistency; what mode reordering
implies for adaptive test-time compute and curated hard subsets; and, if
pre-registered predictions failed, the licensing negative result (reported
ceilings are model properties after all).}

\section*{Reproducibility Statement}
All experiments are inference-only on open-weight models. Tier~A fits in
\textbf{$\approx$12.5 GPU-hours}; Tier~B strengthening adds \textbf{$\approx$16.1 h}
for item expansion, multi-config deep-$N$ (O2), and seed replicates at $c_1$
(O2b). Total paper-critical generation is \textbf{$\approx$28.6 h} ($\approx$32 h
with pilots). Full specification in \texttt{docs/METHOD\_SPEC.md}.

\bibliographystyle{plainnat}
\bibliography{references}

\end{document}
